% ЗАКЛЮЧЕНИЕ, без номер.
% Обобщава решението, предимствата и недостатъците му, резултатите и посоките напред.
\section*{ЗАКЛЮЧЕНИЕ}
\label{sec:conclusion}
\addcontentsline{toc}{section}{ЗАКЛЮЧЕНИЕ}

Проектът изгражда изчислително интензивна услуга, разгърната едновременно при двама
публични облачни доставчика от едно декларативно описание на инфраструктурата, и
измерва разликата между двете среди при контролирано еднакви условия. Разделението
между общата и платформено зависимата част е направено така, че различията между
доставчиците да са локализирани в два модула с еднакъв интерфейс, а всичко останало,
включително контейнерният образ, описанието на работното натоварване и правилото за
мащабиране, да остане едно.

\textbf{Предимства на избраното решение.} Смяната на доставчика е смяна на стойността
на един параметър, което прави сравнението осмислено и повторяемо. Метриките се
събират от самото приложение, с един и същ код при двете разгръщания, така че
разликата в измереното не се дължи на разлика в измервателния уред. В хранилището с
изходния текст не се съхранява нито един секрет.

\textbf{Недостатъци и ограничения.} Изборът на управляван Kubernetes дава най-голяма
преносимост на описанието, но и най-висока постоянна такса, така че изводът за цената
важи за този начин на изпълнение, а не за облака изобщо. Съпоставянето на типове
възли между двамата доставчици е приблизително, защото еднакви типове не съществуват.
Мрежовото разстояние до двата региона не е изолирано напълно.

\TODO{Да се обобщят получените резултати с оглед на целта, поставена в увода, след
като измерванията бъдат проведени. Тук се посочва и кои от предварителните очаквания
не са се потвърдили, ако има такива.}

\textbf{Насоки за по-нататъшна работа.} Естественото продължение е разширяване на
сравнението с трети доставчик, добавяне на второ, входно изходно интензивно
натоварване, при което гаранциите на обектното хранилище влизат в сила по различен
начин, както и проверка на поведението при отказ на цяла зона на достъпност.
